\subsection{Paired comparison with the edit-consistency objective}
\label{sec:paired}

\paragraph{Design, fixed before the run.}
The panel contains the first \SSpanelTrainN{} training groups (\SSpanelTrainOrbits{} orbits) and
\SSpanelDevN{} development groups (\SSpanelDevOrbits{} orbits). The development groups start after
the \SSpanelDevExcluded{} groups already scored in the small-panel stage. The two panels share
\SSpanelOverlap{} orbits, and \PRdevBNN{} of the development groups are binding-necessary. Arm A
uses the hard-negative loss. Arm B adds $\lambda\,\mathcal{L}_{\mathrm{edit}}$ with
$\lambda=\PRlambda$, the smallest non-zero value of the earlier pre-registered grid. We disclose
that this grid was set after observing gradient scales, and $\lambda$ was not changed after any
development result was seen. Both arms share initialization, batch order, supervision,
false-negative masks, cached encoder features, optimizer (AdamW, learning rate \PRlr, weight decay
\PRwd) and \PRsteps{} updates of \PRbatch{} groups. Seeds 0, 1 and 2 are paired optimization
repeats, and the development panel is scored only after the final update. At initialization the
edit gradient was \PRgradInitMean{} times the task gradient (mean over the three first batches), so
$\lambda$ scales it to about \PRgradInitLamMean{} of the task gradient. Encoding the
\PRrenderImages{} panel images and their captions took \PRencCPU{} CPU seconds, and the audit
decoder found \PRrenderMismatch{} mismatches.

\begin{table}[t]
\caption{Paired comparison on \SSpanelTrainN{} training and \SSpanelDevN{} development groups
(\PRdevBNN{} binding-necessary). Seeds 0/1/2 are paired optimization repeats on the same panel.
Chance for independent scores: group $1/6$, GroupMatch $1/2$, pair AUC $0.5$. Source:
\url{runs/paired_v1_20260926T230018Z/paired_summary.json}.}
\label{tab:paired}
\centering\small
% GENERATED by scripts/export_paper_numbers.py from runs/paired_v1_20260926T230018Z; do not edit by hand.
\begin{tabular}{lccc}
\toprule
Seeds 0/1/2 & A: hard-neg. & B: + edit-cons. & B$-$A \\
\midrule
Dev \bn{} group score & .094/.016/.000 & .000/.125/.234 & $-$.094/.109/.234 \\
Dev \bn{} GroupMatch & .688/.641/.656 & .609/.812/.719 & $-$.078/.172/.062 \\
Dev \bn{} pair AUC & .520/.503/.501 & .506/.549/.574 & $-$.014/.045/.073 \\
Dev binding-only group & .156/.000/.000 & .000/.156/.469 & $-$.156/.156/.469 \\
Dev relation-only group & .031/.031/.000 & .000/.094/.000 & $-$.031/.062/.000 \\
Dev all-kinds group & .367/.289/.250 & .297/.383/.461 & $-$.070/.094/.211 \\
Dev all-kinds GroupMatch & .844/.820/.828 & .805/.906/.859 & $-$.039/.086/.031 \\
Train \bn{} group (fit) & .535/.152/.066 & .094/.746/.703 & $-$.441/.594/.637 \\
Final task loss & 0.51/1.83/7.18 & 0.91/0.25/0.29 & 0.40/$-$1.58/$-$6.89 \\
\bottomrule
\end{tabular}

\end{table}

\paragraph{Outcome under the pre-fixed rules.}
Neither arm met the fit criterion, a seed-mean binding-necessary training group score of at least
$0.75$: A reached \PRATrainBNGroupMean{} and B \PRBTrainBNGroupMean{} (Table~\ref{tab:paired}). The
pre-registered verdict is therefore \texttt{\PRverdict}. Training was unstable in both arms. The
task loss oscillated during training, and the largest logged gradient norm per run ranged from
\PRgradMinOfMax{} to \PRgradMaxAll. Arm A finished seed 2 at a task loss of \PRALossLastSeedTwo{},
above its initial value, and seed 1 at \PRALossLastSeedOne. Arm B finished at
\PRBLossLastSeedZero, \PRBLossLastSeedOne{} and \PRBLossLastSeedTwo.

\paragraph{What the difference does and does not show.}
On binding-necessary development groups, the edit-consistency arm scored higher in \PRdiffGroupPos{}
of 3 seeds. The mean paired difference in group score is \PRdiffGroup. Its group-level bootstrap
interval is [\PRdiffGroupLo, \PRdiffGroupHi], but that interval treats the three runs as fixed and
does not reflect their optimization variance. The per-seed differences in Table~\ref{tab:paired} are
positive in the seeds where the baseline ended with a high final loss, and negative in the seed
where it did not. We therefore do not read them as added utility of the objective. Both arms also
sit below the independent-score chance level on the group score (A: \PRACIGroupMean{},
[\PRACIGroupLo, \PRACIGroupHi]; B: \PRBCIGroupMean{}, [\PRBCIGroupLo, \PRBCIGroupHi]) and above it
on GroupMatch (A: \PRACIMatchMean; B: \PRBCIMatchMean). A scorer whose preferences are correlated
across the two captions, for instance one that prefers the same image for both, can fall below
$1/6$ on the group score while still matching better than $1/2$, so the two metrics must be read
together. Zero-shot CLIP on the same subset reaches \PRzsBNGroup{} (group) and \PRzsBNMatch{}
(GroupMatch). Following the pre-fixed rules, the edit-objective branch is put on hold. No
stabilization, encoder change, data increase or new loss was started on the basis of this result.
